\section{仿射开集与局部性质}\label{sec:affine-opens-and-local-properties}

已经证明拟射影簇理论专化为仿射簇理论, 现在可以描述拟射影簇最有用的性质之一: 每个拟射影簇都是\textquotedbl{}局部仿射\textquotedbl{}的, 其确切含义将在本节说明. 这一事实的作用在于, 它使我们能把拟射影簇的许多性质, 例如维数和光滑性, 归结为仿射簇的相应性质. 本书第一部分已经详细建立了这些仿射性质.

证明每个拟射影簇都局部仿射的第一步, 是研究仿射空间的闭子集中的非消没集.

\begin{Proposition}{非消没集是仿射簇}\leavevmode\label{prop:nonvanishing-sets-in-affine-closed-sets-are-affine}
若 $X\subseteq\mathbb{A}^{n}$ 是 Zariski 闭集, 且 $f\in K[x_{1},\ldots,x_{n}]$, 则集合 $U(f)=X\setminus\mathcal{V}(f)$ 是仿射簇.
\end{Proposition}

Example~\ref{ex:punctured-affine-line-isomorphic-to-hyperbola} 展示了此结果的一个实例: 非消没集 $\mathbb{A}^{1}\setminus\mathcal{V}(x)$ 同构于闭集 $\mathcal{V}_{\mathbb{A}}(xy-1)\subseteq\mathbb{A}^{2}$. 更一般地, 类似的论证表明, 对任意 $f\in K[x_{1},\ldots,x_{n}]$, $\mathbb{A}^{n}\setminus\mathcal{V}(f)$ 同构于闭集 $\mathcal{V}(x_{n+1}f-1)\subseteq\mathbb{A}^{n+1}$. 这些特例为一般证明提供了范例, 现在给出证明.

\begingroup
\renewcommand{\AJProofBegin}{%
  \par\addvspace{\topsep}\noindent
  {\sffamily\bfseries\CJKfamily{pffont}Proof of Proposition~\ref{prop:nonvanishing-sets-in-affine-closed-sets-are-affine}}\enspace\ignorespaces
}
\begin{Proof}
令 $X\subseteq\mathbb{A}^{n}$ 为闭集, 从而可以写成
\[
  X=\mathcal{V}_{\mathbb{A}}(g_{1},\ldots,g_{m}),
\]
其中 $g_{1},\ldots,g_{m}\in K[x_{1},\ldots,x_{n}]$. 给定 $f\in K[x_{1},\ldots,x_{n}]$, 定义闭集 $Y\subseteq\mathbb{A}^{n+1}$ 为
\[
  Y=\mathcal{V}_{\mathbb{A}}(g_{1},\ldots,g_{m},x_{n+1}f-1).
\]
我们断言 $U(f)$ 与 $Y$ 作为拟射影簇同构. 由于 $Y$ 是仿射空间的闭子集, 这将蕴含 $U(f)$ 是仿射簇.

要看出 $U(f)\cong Y$, 考察投影映射
\[
  \begin{aligned}
    H:Y&\to U(f)\\
    (a_{1},\ldots,a_{n},a_{n+1})&\mapsto(a_{1},\ldots,a_{n}).
  \end{aligned}
\]
注意 (Exercise~\ref{ex:graph-of-inverse-polynomial-projection-is-bijective}), $H$ 为双射, 其逆 $H^{-1}$ 由下式给出:
\[
  (a_{1},\ldots,a_{n})\mapsto\left(a_{1},\ldots,a_{n},\frac{1}{f(a)}\right).
\]
只需再论证 $H$ 和 $H^{-1}$ 都正则. 为此, 将 $Y$ 和 $U(f)$ 分别与它们在 $\mathbb{P}^{n+1}$ 和 $\mathbb{P}^{n}$ 中的像等同. 令 $\overline{f}$ 为 $f$ 的齐次化, 并设 $\deg(\overline{f})=d$. 则 $H$ 对应于正则映射
\[
  [a_{0}:\cdots:a_{n}:a_{n+1}]\mapsto[a_{0}:\cdots:a_{n}],
\]
而稍加思考可知 $H^{-1}$ 对应于正则映射
\[
  [a_{0}:\cdots:a_{n}]\mapsto[a_{0}\overline{f}(a):\cdots:a_{n}\overline{f}(a):a_{0}^{d+1}].
\]
因此 $U(f)\cong Y$, 如所断言.
\end{Proof}
\endgroup

\BookPairedNoteLeft{
上述证明的关键在于, 转到射影空间后, 可以消去分母中的 $f(a)$.
}{
Proposition~\ref{prop:nonvanishing-sets-in-affine-closed-sets-are-affine} 表明, 仿射空间闭子集中的非消没集本身是仿射簇. 这并不对仿射空间的所有开子集成立. 例如, Example~\ref{ex:punctured-affine-plane-is-not-affine} 中遇到的拟射影簇 $\mathbb{A}^{2}\setminus\{(0,0)\}$ 在 $\mathbb{A}^{2}$ 中是开的, 但不是仿射簇.
}
现在利用 Proposition~\ref{prop:nonvanishing-sets-in-affine-closed-sets-are-affine} 证明本节的主要结果: 每个拟射影簇都局部同构于一个仿射簇, 即拟射影簇的每一点都包含在一个仿射开子簇中. 事实上, 我们证明更强的结果: 每个拟射影簇都能由有限多个仿射开子簇覆盖.

\begin{Theorem}{拟射影簇是局部仿射的}\leavevmode\label{thm:quasiprojective-varieties-have-finite-affine-open-covers}
设 $X$ 为拟射影簇. 则存在有限多个仿射开子簇 $U_{1},\ldots,U_{k}\subseteq X$, 使 $X=U_{1}\cup\cdots\cup U_{k}$. 特别地, 对每个 $p\in X$, 都存在仿射开子簇 $U\subseteq X$ 使 $p\in U$.
\end{Theorem}
\begin{Proof}
令 $X\subseteq\mathbb{P}^{n}$ 为拟射影簇. 按定义, 有
\[
  X=Y\cap(\mathbb{P}^{n}\setminus Z)=Y\setminus Z,
\]
其中 $Y,Z\subseteq\mathbb{P}^{n}$ 为射影簇. 对某些齐次多项式 $f_{1},\ldots,f_{k}\in K[x_{0},\ldots,x_{n}]$, 写成 $Z=\mathcal{V}(f_{1},\ldots,f_{k})$. 对每个 $i=1,\ldots,k$ 及每个 $j=0,\ldots,n$, 定义
\[
  U_{i,j}=(Y\setminus\mathcal{V}(f_{i}))\cap\mathbb{A}^{n}_{j}.
\]
注意 $Z=\mathcal{V}(f_{1},\ldots,f_{k})\subseteq\mathcal{V}(f_{i})$, 所以 $Y\setminus\mathcal{V}(f_{i})=X\setminus\mathcal{V}(f_{i})$. 因此, 对每个 $i$, $Y\setminus\mathcal{V}(f_{i})$ 是 $X$ 的开子簇. 又因为 $\mathbb{A}^{n}_{j}$ 在 $\mathbb{P}^{n}$ 中开, 对所有 $i,j$, $U_{i,j}$ 都是 $X$ 的开子簇. 此外, 注意
\[
  U_{i,j}=(Y\cap\mathbb{A}^{n}_{j})\setminus\mathcal{V}(f_{i,j}),
\]
其中 $f_{i,j}$ 是在 $f_{i}$ 中令 $x_{j}=1$ 得到的 $n$ 元多项式. 由于 $Y\cap\mathbb{A}^{n}_{j}$ 是 $\mathbb{A}^{n}_{j}$ 的闭子集, 由 Proposition~\ref{prop:nonvanishing-sets-in-affine-closed-sets-are-affine} 可知每个 $U_{i,j}$ 都是仿射簇. 定理现在由下式得到:
\[
  X=\bigcup_{\substack{1\le i\le k\\0\le j\le n}}U_{i,j},
\]
建议读者在 Exercise~\ref{ex:finite-affine-open-cover-from-defining-equations} 中验证这一等式, 从而证明 $X$ 的每一点都属于 (至少) 一个仿射开子簇 $U_{i,j}$.
\end{Proof}

回顾第~\ref{sec:quasiprojective-varieties}~节, 点 $p\in X$ 的\emph{邻域}是包含 $p$ 的 $X$ 的任一开子集. 因此, Theorem~\ref{thm:quasiprojective-varieties-have-finite-affine-open-covers} 可以重述为: \emph{拟射影簇的每一点都有一个仿射邻域}. 下面的例子具体说明这一定理.

\begin{Example}{$\mathbb{A}^{2}\setminus\{(0,0)\}$ 是局部仿射的}\leavevmode\label{ex:punctured-affine-plane-has-affine-open-cover}
令 $X=\mathbb{A}^{2}\setminus\{(0,0)\}$, 以通常的方式将它视为 $\mathbb{P}^{2}$ 中的拟射影簇. 由 Proposition~\ref{prop:nonvanishing-sets-in-affine-closed-sets-are-affine}, 子集
\[
  U=X\setminus\mathcal{V}(x)=\mathbb{A}^{2}\setminus\mathcal{V}(x)
\]
是仿射簇. 特别地, $U\cong\mathcal{V}(xz-1)\subseteq\mathbb{A}^{3}$, 同构由
\[
  (x,y)\mapsto\left(x,y,\frac{1}{x}\right)
\]
给出. 类似地,
\[
  V=X\setminus\mathcal{V}(y)=\mathbb{A}^{2}\setminus\mathcal{V}(y)
\]
也是仿射簇. $U$ 和 $V$ 都是 $X$ 的 Zariski 开子簇. 由于 $X$ 的每一点的 $x$ 坐标或 $y$ 坐标必有一个非零, 每一点都包含在 $U$ 或 $V$ 中. 这样便验证了拟射影簇 $\mathbb{A}^{2}\setminus\{(0,0)\}$ 的每一点都包含在一个仿射邻域中.
\end{Example}

如本节开头所说, Theorem~\ref{thm:quasiprojective-varieties-have-finite-affine-open-covers} 的作用在于, 它允许把拟射影簇某些性质的研究归结到仿射情形. 现在为这些性质命名.

\begin{Definition}{局部性质}\leavevmode\label{def:local-property-of-topological-spaces}
令 $\mathcal{C}$ 为一类拓扑空间. 若 $\mathcal{C}$ 上的性质 $\mathcal{P}$ 满足以下条件, 则称其为\emph{局部的}: 若 $X\in\mathcal{C}$, 且每个 $p\in X$ 都包含在具有性质 $\mathcal{P}$ 的开集 $U_{p}\subseteq X$ 中, 则 $X$ 具有性质 $\mathcal{P}$.
\end{Definition}

一个反例有助于理解这一定义: 仿射性不是拟射影簇类上的局部性质, 因为例如 $\mathbb{P}^{n}$ 不是仿射簇 (Exercise~\ref{ex:positive-dimensional-projective-space-is-not-affine}), 但 $\mathbb{P}^{n}$ 的每一点都包含在某个仿射开片中. 打个比方, 如果一个人的视力很差, 看不到任何给定点的一个小邻域之外的地方, 就不可能辨认自己站在仿射簇 $\mathbb{A}^{n}$ 中还是非仿射簇 $\mathbb{P}^{n}$ 中. 这有点像早期人类无法看出自己站的不是平坦地面构成的地球, 而是一个 (近似) 球形的地球.

另一方面, 光滑性是不可约拟射影簇类上的局部性质的一个例子, 如我们将在下一章看到的. 直观上, 若 $X$ 有一个奇点, 那么即使是 $X$ 中一个近视的居民, 只要站在这个点上, 也会知道它是奇点. 类似地, 具有某个特定维数是不可约簇的局部性质: 在不可约簇内可独立移动的方向数, 可以在任意一点的邻域中判定. 下一章将深入讨论这两个性质.

本节最后给出局部性质的两个有用结果, 下一节都会用到. 第一个引理说, 要验证拓扑空间的一个子集是闭的, 只须检查它在每一点的某个邻域中是闭的.

\begin{Lemma}{闭性是局部的}\leavevmode\label{lem:closedness-is-local}
设 $Y$ 为拓扑空间, $X\subseteq Y$ 为子集. 若每个 $p\in Y$ 都包含在一个开集 $U_{p}\subseteq Y$ 中, 使 $X\cap U_{p}$ 在 $U_{p}$ 中闭, 则 $X$ 在 $Y$ 中闭.
\end{Lemma}
\begin{Proof}
对每个 $p\in Y$, 固定开集 $U_{p}$, 使得 $p\in U_{p}$ 且 $X\cap U_{p}$ 在 $U_{p}$ 中闭. 由开集的定义, 对每个 $p$, 存在闭集 $Z_{p}\subseteq Y$, 使 $U_{p}=Y\setminus Z_{p}$. 由子空间拓扑的定义, 存在闭集 $Z'_{p}\subseteq Y$, 使 $X\cap U_{p}=Z'_{p}\cap U_{p}$. 现在断言
\[
  X=\bigcap_{p\in Y}(Z_{p}\cup Z'_{p}),
\]
由此可知 $X$ 在 $Y$ 中闭.

($\subseteq$) 假设 $q\in X$. 则对每个 $p\in Y$, 有 $q\in U_{p}$ 或 $q\notin U_{p}$. 若 $q\in U_{p}$, 则 $q\in X\cap U_{p}=Z'_{p}\cap U_{p}$, 故 $q\in Z'_{p}$. 若 $q\notin U_{p}$, 则 $q\in Z_{p}$. 因此无论哪种情形, 都有 $q\in Z_{p}\cup Z'_{p}$, 且这对所有 $p\in Y$ 成立.

($\supseteq$) 反过来, 假设对所有 $p\in Y$ 都有 $q\in Z_{p}\cup Z'_{p}$. 由于 $q\in U_{q}$, 有 $q\notin Z_{q}$. 因而 $q\in Z_{q}\cup Z'_{q}$ 蕴含 $q\in Z'_{q}$. 于是 $q\in Z'_{q}\cap U_{q}=X\cap U_{q}$, 所以 $q\in X$.
\end{Proof}

下一个引理断言, 正则性可以在定义域上局部检查.

\begin{Lemma}{正则性在定义域上是局部的}\leavevmode\label{lem:regularity-is-local-on-domain}
令 $X\subseteq\mathbb{P}^{m}$ 和 $Y\subseteq\mathbb{P}^{n}$ 为拟射影簇, $F:X\to Y$ 为映射. 若每个 $p\in X$ 都包含在一个开集 $U_{p}\subseteq X$ 中, 使限制 $F:U_{p}\to Y$ 正则, 则 $F$ 正则.
\end{Lemma}
\begin{Proof}
设映射 $F:X\to Y$ 满足引理的假设. 为证明 $F$ 正则, 令 $p\in X$. 按假设, 存在开集 $U_{p}\subseteq X$ 和同次数的齐次多项式 $g_{0},\ldots,g_{n}\in K[x_{0},\ldots,x_{m}]$, 使 $p\notin\mathcal{V}(g_{0},\ldots,g_{n})$, 且
\[
  F(a)=[g_{0}(a):\cdots:g_{n}(a)]
  \qquad\text{对所有 }a\in U_{p}\setminus\mathcal{V}(g_{0},\ldots,g_{n}).
\]
要使 $F$ 满足正则映射的定义, 上式须在整个 $X\setminus\mathcal{V}(g_{0},\ldots,g_{n})$ 上成立, 但它未必如此.

为解决这一问题, 选取齐次多项式 $f\in K[x_{0},\ldots,x_{n}]$, 使 $f$ 在闭集 $X\setminus U_{p}\subseteq X$ 上消没, 且 $f(p)\ne0$. 对每个 $i=0,\ldots,n$, 定义 $f_{i}=fg_{i}$. 注意两个函数 $[f_{0}:\cdots:f_{n}]$ 和 $[g_{0}:\cdots:g_{n}]$ 在它们都有定义之处取值相同, 但前者的定义域较小. 特别地, 由于 $f$ 在 $X\setminus U_{p}$ 上消没, 得
\[
  X\setminus\mathcal{V}(f_{0},\ldots,f_{n})
    =U_{p}\setminus\mathcal{V}(f_{0},\ldots,f_{n}).
\]
因此, 由构造可知 $p\notin X\setminus\mathcal{V}(f_{0},\ldots,f_{n})$, 且
\[
  F(a)=[f_{0}(a):\cdots:f_{n}(a)]
  \qquad\text{对所有 }a\in X\setminus\mathcal{V}(f_{0},\ldots,f_{n}),
\]
所以 $F$ 在 $p$ 处正则.
\end{Proof}

\BookExerciseHeading
\begin{BookExercises}
  \item\label{ex:graph-of-inverse-polynomial-projection-is-bijective} 令 $g_{1},\ldots,g_{m},f\in K[x_{0},\ldots,x_{n}]$, 定义
  \[
    \begin{aligned}
    X&=\mathcal{V}(g_{1},\ldots,g_{m})\subseteq\mathbb{A}^{n},\\
    Y&=\mathcal{V}(g_{1},\ldots,g_{m},x_{n+1}f-1)\subseteq\mathbb{A}^{n+1}.
    \end{aligned}
  \]
  证明函数
  \[
    \begin{aligned}
      \phi:Y&\to X\setminus\mathcal{V}(f)\\
      (a_{1},\ldots,a_{n},a_{n+1})&\mapsto(a_{1},\ldots,a_{n})
    \end{aligned}
  \]
  是双射, 并计算它的逆.
  \item\label{ex:regular-function-ring-of-principal-open-is-localization} 令 $X\subseteq\mathbb{A}^{n}$ 为闭集, $f\in K[x_{1},\ldots,x_{n}]$. 置 $F=[f]\in K[X]$, 证明
  \[
    K[X\setminus\mathcal{V}(f)]=\frac{K[X][z]}{\langle zF-1\rangle}.
  \]
  (当 $a$ 在 $R$ 中不是幂零元时, 通常用 $R[a^{-1}]$ 记环 $R[z]/\langle za-1\rangle$. 因此, 有 $K[X\setminus\mathcal{V}(f)]=K[X][F^{-1}]$.)
  \item\label{ex:finite-affine-open-cover-from-defining-equations} 设 $Y,Z\subseteq\mathbb{P}^{n}$ 为射影簇, 且 $Z=\mathcal{V}(f_{1},\ldots,f_{k})$. 对每个 $i=1,\ldots,k$ 和 $j=0,\ldots,n$, 定义
  \[
    U_{i,j}=(Y\setminus\mathcal{V}(f_{i}))\cap\mathbb{A}^{n}_{j}.
  \]
  证明
  \[
    Y\setminus Z=\bigcup_{\substack{1\le i\le k\\0\le j\le n}}U_{i,j}.
  \]
  \item\label{ex:affine-open-cover-of-punctured-affine-space} 构造 $\mathbb{A}^{n}\setminus\{(0,\ldots,0)\}$ 的有限仿射开覆盖.
  \item\label{ex:affine-open-cover-of-open-projective-quadric} 令 $Y=\mathcal{V}(x^{2}+y^{2}-zw)\subseteq\mathbb{P}^{3}$, $Z=\mathcal{V}(x,y,z)\subseteq\mathbb{P}^{3}$. 构造 $X=Y\setminus Z$ 的有限仿射开覆盖.
  \item\label{ex:projective-hypersurface-complement-is-affine} 令 $f\in K[x_{0},\ldots,x_{n}]$ 为非常数齐次多项式. 证明 $\mathbb{P}^{n}\setminus\mathcal{V}(f)$ 是仿射簇. (提示: 使用一个 Veronese 映射.)
\end{BookExercises}
